% Generated draft: requires mathematical and source review.
\documentclass[11pt]{article}
\usepackage[margin=25mm]{geometry}
\usepackage{fontspec}
\setmainfont{Latin Modern Roman}
\usepackage{amsmath,amssymb,mathtools,unicode-math}
\setmathfont{Latin Modern Math}
\usepackage{enumitem,needspace,xurl,hyperref}
\setlength{\parindent}{0pt}
\setlength{\parskip}{5pt}
\setlength{\emergencystretch}{3em}
\newcommand{\N}{\mathbb{N}}
\newcommand{\Z}{\mathbb{Z}}
\newcommand{\Q}{\mathbb{Q}}
\newcommand{\R}{\mathbb{R}}
\newcommand{\C}{\mathbb{C}}
\newcommand{\F}{\mathbb{F}}
\newcommand{\A}{\mathbb{A}}
\newcommand{\PP}{\mathbb{P}}
\newcommand{\E}{\mathbb{E}}
\newcommand{\eps}{\varepsilon}
\newcommand{\dd}{\,\mathrm{d}}
\begin{document}
\section*{UnsolvedMath ID: 30004134}
\textbf{Source problem code:} catalogue entry
\subsection{Definitions and mathematical statement}
Fix integers \(k\ge r\ge2\) satisfying \(r\ge3\) or \(k\ge4\). Under the conventional finite simple hypergraph interpretation, let \(H=(V,E)\), where \(V\) is finite, \(E\subseteq\binom{V}{k}\), and \(m=|E|\). For a partition \(V=V_1\sqcup\cdots\sqcup V_r\), define its \(r\)-cut to consist of the edges meeting every part, with size
\[
C_H(V_1,\ldots,V_r)=\bigl|\{e\in E:e\cap V_i\ne\varnothing\text{ for every }i\in\{1,\ldots,r\}\}\bigr|.
\]
Here parts may be empty; any cut of positive size necessarily has all parts nonempty. Let \(S(k,r)\) denote the Stirling number of the second kind, the number of partitions of a \(k\)-element set into \(r\) nonempty unlabelled blocks, and set
\[
\alpha_{r,k}=\frac{S(k,r)r!}{r^k}.
\]
Thus \(\alpha_{r,k}\) is the probability that an edge meets all \(r\) classes when its vertices receive independent uniform labels from \(\{1,\ldots,r\}\).

The question is whether, for every fixed admissible pair \((r,k)\), there exist a constant \(c_{r,k}>0\) and an integer \(m_0(r,k)\ge1\) such that every such hypergraph with \(m\ge m_0(r,k)\) admits a partition satisfying
\[
C_H(V_1,\ldots,V_r)\ge\alpha_{r,k}m+c_{r,k}m^{2/3}.
\]
The constants must be independent of \(H\), \(m\), and \(|V|\). This is the asymptotic interpretation of the supplied term \(\Omega(m^{2/3})\).

The record does not explicitly define an \(r\)-cut or specify finiteness and the exclusion of repeated edges. The definitions above give a conventional interpretation requiring source confirmation; they do not establish that these conventions are explicitly stated in the cited report. The supplied label \texttt{partially\_solved} and its literature assessment are dated curation, not mathematical conclusions verified by this draft.
\subsection{Short English statement}
For each fixed admissible \((r,k)\), must every sufficiently large \(k\)-uniform hypergraph admit a partition into \(r\) classes for which the number of edges meeting every class exceeds the random-partition baseline \(\alpha_{r,k}m\) by at least a positive constant times \(m^{2/3}\)? Hypergraph and cut conventions require confirmation from the source.
\subsection{Sources}
\begin{itemize}
\item Supplied local UnsolvedMath registry record, permanent ID \(30004134\), source problem code \texttt{catalogue entry}, title \emph{Second-Order Cut Bounds for Uniform Hypergraphs}; collection: Oberwolfach Reports Open Problems. Local registry paths: \texttt{data/unsolvedmath/problems.json} and the supplementary \texttt{catalog/problems.json}. The supplied proposed year is \(2019\); no individual proposer is named.
\item Supplied report citation: \emph{Combinatorics, Probability and Computing} (2020), Oberwolfach Reports, DOI 10.4171/owr/2019/19, pp.\ 1113--1189. \url{https://doi.org/10.4171/owr/2019/19}. The background describes these reports as documenting mathematical workshops and their open-problem discussions. This supplied citation has not been independently consulted.
\item Supplied literature source: Conlon--Fox--Kwan--Sudakov, \url{https://arxiv.org/abs/1803.08462}.
\item Supplied literature source: Janzer--Portier, \url{https://arxiv.org/abs/2511.08501}.
\item Literature assessment dated exactly \texttt{2026-08-22}: the supplied curation labels the problem partially solved and states that the exact conjectured \(m^{2/3}\) bound is not known in general, that special cases have stronger bounds, and that recent work proves an approximate \(m^{2/3-\varepsilon}\) result in broad ranges. This is supplied dated curation derived from the OWR record, not independently verified; its time-sensitive status and cited claims require independent checking.
\item The supplied statement status is \texttt{corrected\_verified}; its accompanying note says mathematical content was preserved while notation, typography, and sentence structure were repaired for publication. This is a source-record assertion, not independent mathematical verification by this draft.
\end{itemize}
\end{document}
